% Auto-generated Beamer translation of a PDF slide deck (Week 1)
% Figures replaced with placeholders as requested.
\documentclass[aspectratio=169]{beamer}

\usetheme{default}
\usecolortheme{default}
\usepackage[T1]{fontenc}
\usepackage{lmodern}
\usepackage{amsmath,amssymb}
\usepackage{hyperref}
\usepackage{graphicx}

% Simple figure placeholder macro
\newcommand{\figplaceholder}[2][0.75\linewidth]{%
  \begin{center}
    \setlength{\fboxsep}{6pt}%
    \fbox{\parbox[c][0.32\textheight][c]{#1}{\centering\small\textit{Figure placeholder:} #2}}
  \end{center}
}

\title{Introduction to Quantum Computing}
\author{MHJ}
\date{}

\begin{document}

% --- Slide 4 ---
\begin{frame}{What is quantum computing?}
\begin{itemize}
  \item It is computation (and more generally, information processing) based on the principles of quantum mechanics, rather than classical physics.
  \item Quantum mechanics is a description of nature at its most fundamental level.
  \begin{itemize}
    \item Formulated in the early 20th century to explain the behavior of subatomic particles.
    \item QM, and its specializations (quantum field theory, quantum chromodynamics, etc.), have been spectacularly successful in explaining microscopic physical phenomena.
  \end{itemize}
\end{itemize}
\figplaceholder{Hydrogen atom; Superconductivity; Bose--Einstein condensates (images on the right)}
\end{frame}

% --- Slide 5 ---
\begin{frame}{The exponentiality of QM}
\begin{itemize}
  \item After nearly a century of study, the best (classical) methods for predicting the behavior of general quantum systems require exponential resources.
  \item The state of $N$ particles requires at least $2^{N}$ numbers to describe.
  \item For $N \sim 300$ (the number of particles in a single uranium atom),
  \[
    2^{300} \gg \approx 10^{82}\quad (\text{\# of atoms in the observable universe}).
  \]
  \item Nature seems to be doing extravagant amounts of computation “behind the scenes”!
\end{itemize}
\end{frame}

% --- Slide 6 ---
\begin{frame}{The exponentiality of QM}
\begin{itemize}
  \item So how do physicists actually do (quantum) physics?
  \item Answer \#1: 100+ years of extremely clever approximations and shortcuts for calculations.
  \begin{itemize}
    \item Bethe ansatz, Feynman diagrams, perturbation theory, mean field approximations, \dots
  \end{itemize}
  \item Answer \#2: computer simulations
  \begin{itemize}
    \item Density functional theory, Quantum Markov Chain Monte Carlo, \dots
  \end{itemize}
  \item Answer \#3: studying systems that allow for Answers \#1 and \#2 above
  \begin{itemize}
    \item Systems that are mostly classical
    \item Systems with non-strongly interacting particles
    \item Lucky that most things we’ve studied so far fall in this category!
  \end{itemize}
  \item But simulating \emph{general} quantum physics still seems like a hard task\dots
\end{itemize}
\figplaceholder{Illustrations on right (plots + cartoon)}
\end{frame}

% --- Slide 7 ---
\begin{frame}{The exponentiality of QM}
\begin{itemize}
  \item In 1981, in his talk titled \emph{“Simulating Physics with Computers”}, Richard Feynman asked the following question:
  \begin{quote}
    “Can probabilistic computers simulate quantum mechanics?”
  \end{quote}
  \item His conclusion, after a lengthy exploration:
  \begin{quote}
    “\dots quantum mechanics [doesn’t] seem to be imitable by a local classical computer.”
  \end{quote}
  \item Instead we need a “computer\dots built of quantum mechanical elements which obey quantum mechanical laws.”
  \item Feynman, along with David Deutsch, Paul Benioff and Yuri Manin, are credited with the idea of computing based on quantum mechanical principles.
\end{itemize}
\end{frame}

% --- Slide 8 ---
\begin{frame}{The early days}
\begin{itemize}
  \item Feynman’s motivation for quantum computers is the most obvious one: simulating quantum systems (almost tautological!).
  \item To physicists in the '80s, this idea might have seemed obvious. To a computer scientist in the '80s, QCs would have seemed like a wild crackpot idea.
  \item Also in the '80s, David Deutsch contemplated other applications of quantum computers.
  \begin{itemize}
    \item He gave an example of a (classical) problem where the “parallel universes” of QM seem to speed up (by a constant factor).
  \end{itemize}
  \item 1992: Dan Simon came up with an example of a problem that could be solved exponentially faster on a QC.
\end{itemize}
\end{frame}

% --- Slide 9 ---
\begin{frame}{Shor’s breakthrough}
\begin{itemize}
  \item 1993: Peter Shor realized that Simon’s idea could be extended to efficiently solve the following problem:
  \item \textbf{(Integer Factorization)} Given integer $N>0$, find integer $M>1$ such that $M$ divides $N$.
  \item The security of the RSA cryptosystem (the most widely deployed public-key encryption on the internet) crucially relies on the assumption that factoring integers is hard!
\end{itemize}
\figplaceholder{Photo of Peter Shor + PayPal/RSA illustration}
\end{frame}

% --- Slide 10 ---
\begin{frame}{Crossroads}
Since Shor’s algorithm, physicists and computer scientists have been faced with three options:
\begin{enumerate}
  \item Quantum mechanics is wrong.
  \item There is a fast classical algorithm for factoring.
  \item Quantum computers are more powerful than classical computers.
\end{enumerate}
\bigskip
\begin{center}
\large At least one of these must be true!\\
\medskip
Which do you think is most likely to be true?
\end{center}
\figplaceholder{Crossroads sign image (bottom-right)}
\end{frame}

% --- Slide 11 ---
\begin{frame}{Crossroads}
\begin{itemize}
  \item Option \#1: QM is wrong
  \begin{itemize}
    \item Perhaps the most successful theory of nature to date.
    \item In all its domains of applicability, we have never found experimental disagreement.
  \end{itemize}
  \item Option \#2: Polynomial-time classical factoring algorithm.
  \begin{itemize}
    \item RSA has been out for nearly 50 years. There is enormous economic pressure to discover weaknesses, if it exists.
    \item The “Gauss argument”: if the Prince of Mathematics couldn’t find it, then it probably doesn’t exist.
  \end{itemize}
\end{itemize}
\end{frame}

% --- Slide 12 ---
\begin{frame}{Crossroads}
\begin{itemize}
  \item Option \#3: Quantum computers are more powerful than classical computers.
  \item This would refute the Extended Church--Turing Thesis.
  \item Church--Turing Thesis: A Turing machine can simulate any effective computation process.
  \item Extended Church--Turing Thesis: A probabilistic Turing machine can efficiently simulate any effective computation process.
  \item Option \#3 would violate the ECT (assuming factoring is hard for classical computers)!
\end{itemize}
\end{frame}

% --- Slide 13 (reprise) ---
\begin{frame}{Crossroads}
Since Shor’s algorithm, physicists and computer scientists have been faced with three options:
\begin{enumerate}
  \item Quantum mechanics is wrong.
  \item There is a fast classical algorithm for factoring.
  \item Quantum computers are more powerful than classical computers.
\end{enumerate}
\bigskip
\begin{center}
\large At least one of these must be true!
\end{center}
\figplaceholder{Crossroads sign image (bottom-right)}
\end{frame}

% --- Slide 14 ---
\begin{frame}{Present day}
Quantum computing is extremely active and exciting:
\begin{enumerate}
  \item Engineering: primitive quantum computers are being built.
  \item Theoretical developments: new algorithms, new cryptographic protocols, new insights into how quantum computing compares to classical computing.
  \item Deep connections to other sciences: condensed matter physics, quantum gravity, materials science, chemistry, computer science, pure mathematics.
\end{enumerate}
\end{frame}

% --- Slide 15 ---
\begin{frame}{Emerging quantum computers}
\begin{itemize}
  \item Many players racing to build (scalable) quantum computers.
  \item Different labs/companies are betting on different horses.
  \item Superconducting qubits, ion traps, photonic systems, topological qubits, \dots
\end{itemize}
\figplaceholder{Company logos / ecosystem figure}
\end{frame}

% --- Slide 16 ---
\begin{frame}{Emerging quantum computers}
\begin{itemize}
  \item Until 2017 or so, most quantum computing devices had $\sim 10$ qubits.
  \item Suddenly, there was a flood of resources devoted to engineering efforts, and now we are seeing programmable quantum computers with $50+$ qubits.
  \item Major problem: devices are insanely noisy!
  \item We are in the “NISQ” era:
  \begin{itemize}
    \item Noisy Intermediate-Scale Quantum
    \item Noisy devices with $\sim 100$ qubits.
  \end{itemize}
  \item Not capable of running Shor’s algorithm, but should still be capable of solving interesting, hard problems.
\end{itemize}
\end{frame}

% --- Slide 17 ---
\begin{frame}{Quantum supremacy}
\begin{itemize}
  \item In Fall 2019, Google announced that they had achieved “Quantum Supremacy” using their 53-qubit Sycamore quantum processor.
  \item \textbf{Quantum supremacy:} a moment when there is a convincing real-world demonstration of a quantum computing task that cannot feasibly be performed by a classical computer.
  \item Sweet spot: 53 qubits just beyond the reach of Google’s compute clusters to simulate in a reasonable time. ($\sim 2^{53}$ ops)
\end{itemize}
\figplaceholder{Nature cover + hardware photos}
\end{frame}

% --- Slide 18 ---
\begin{frame}{Summary of hardware efforts}
\begin{itemize}
  \item We’re in the NISQ era, but we have compelling evidence that we can already perform super-classical tasks using these machines.
  \item There’s a “qubit race” going on, but qubit count is not everything!
  \item Scaling up (i.e. more, higher-quality qubits) is a tough engineering challenge, but no fundamental obstacles to building a QC with $10^{6}$ noiseless qubits.
  \item Still, large-scale fault-tolerant QCs look like they’re many years away.
\end{itemize}
\bigskip
\begin{center}
\large What interesting problems can we solve on near-term quantum computers?
\end{center}
\figplaceholder{Hardware-related images on right + green callout}
\end{frame}

% --- Slide 19 ---
\begin{frame}{Quantum algorithms}
\begin{itemize}
  \item Exponential speedups for structured, algebraic problems
  \begin{itemize}
    \item Factoring (Shor’s algorithm)
    \item Hidden subgroup problem
  \end{itemize}
  \item Polynomial speedups for unstructured search problems
  \begin{itemize}
    \item Grover search
  \end{itemize}
  \item Exponential speedups for quantum simulation
  \begin{itemize}
    \item Simulating quantum physics and chemistry (Feynman’s dream)
    \item Probably the most important application of quantum computers so far!
  \end{itemize}
\end{itemize}
\figplaceholder{Right-side diagrams; includes $N=pq$ and circuit/chemistry graphics}
\end{frame}

% --- Slide 20 ---
\begin{frame}{Quantum algorithms}
\begin{itemize}
  \item Algorithms to be run on near-term QCs:
  \begin{itemize}
    \item Variational eigensolvers
    \item Classical--quantum hybrid algorithms
    \item Quantum Machine Learning
    \item Linear system solvers / SDPs / Convex Optimization
    \item Quantum neural nets
    \item Recommendation systems
    \item \dots more
  \end{itemize}
\end{itemize}
\bigskip
\begin{center}
\large What types of problems admit a quantum advantage?
\end{center}
\end{frame}

% --- Slide 21 ---
\begin{frame}{Connections with fundamental physics}
\begin{itemize}
  \item The key to a theory of Quantum Gravity could be quantum entanglement, and quantum error correcting codes.
  \item The “Blackhole Firewall Paradox” is an issue about quantum information, and possible resolutions involve quantum cryptography.
\end{itemize}
\bigskip
\begin{center}
\large What can Quantum Computing tell us about Nature?
\end{center}
\figplaceholder{Cylinder/geometry graphic + astronaut/space image}
\end{frame}

% --- Slide 22 ---
\begin{frame}{Uncharted territory}
\figplaceholder[0.9\linewidth]{Large landscape photo}
\end{frame}

% --- Slide 23 ---
\begin{frame}{What you’ll get out of this class}
\begin{itemize}
  \item Learn the basic principles of quantum information and computation
  \item Learn fundamental quantum algorithms and quantum protocols
  \item Get an idea of the current state of the field, and
  \item Get an idea of what the big questions are.
\end{itemize}
\end{frame}

% --- Slide 24 ---
\begin{frame}{Topics covered}
\begin{itemize}
  \item Basics of quantum information
  \item Quantum circuit model
  \item Essential algorithms
  \begin{itemize}
    \item Shor, Grover, Quantum Fourier Transform
  \end{itemize}
  \item Entanglement and its uses
  \item Quantum error-correction
  \item (Time permitting) Quantum complexity theory, quantum cryptography
\end{itemize}
\end{frame}

% --- Slide 25 ---
\begin{frame}{Target audience}
\textbf{Who this class is for:}
\begin{itemize}
  \item Computer scientists, physicists, mathematicians, chemists, \dots
  \item People with a solid linear algebra and probability background
  \item People who want to learn about the theory of quantum computing and quantum information.
\end{itemize}
\medskip
\textbf{Who this class is NOT for:}
\begin{itemize}
  \item People who want to learn how to build a quantum computer
  \item People who want to learn quantum physics
\end{itemize}
\end{frame}


% --- Slide 33 ---
\begin{frame}{Starting point}
\begin{center}
\large
Quantum information theory is a generalization of classical probability theory where probabilities can be negative, or even complex numbers.
\end{center}
\end{frame}

% --- Slide 34 ---
\begin{frame}{Starting point}
\begin{itemize}
  \item Consider a physical system with $d$ distinguishable states, numbered $0,1,\dots,d-1$.
  \item There is also an observer external to the system.
  \item There are two things that can occur:
  \begin{itemize}
    \item \textbf{Measurement:} the external observer can measure the state of $S$.
    \item \textbf{Isolated evolution:} the system can change, without interacting with the external observer.
  \end{itemize}
\end{itemize}
\end{frame}

% --- Slide 35 ---
\begin{frame}{Classical physics}
\begin{itemize}
  \item Initially, the observer assigns a state to the system $S$.
  \item According to classical physics, we can model the state of $S$ as a probability distribution over states, represented as a column vector:
  \[
    \mathbf{p}=
    \begin{pmatrix}
      p_0\\ \vdots\\ p_{d-1}
    \end{pmatrix}\in \mathbb{R}^d,
    \qquad p_i\ge 0,\quad \sum_{i=0}^{d-1}p_i=1.
  \]
\end{itemize}
\end{frame}

% --- Slide 36 ---
\begin{frame}{Classical physics}
\begin{itemize}
  \item If the observer measures (i.e.\ “take a peek at”) the system $S$, then obtains a measurement outcome $i$ with probability $p_i$,
  \item and then the state of the system gets updated to
  \[
    \mathbf{p}\mapsto \mathbf{p}' = \mathbf{e}_i
    =
    \begin{pmatrix}
      0\\ \vdots\\ 1\\ \vdots\\ 0
    \end{pmatrix},
  \]
  where the $1$ is in the $i$th position.
  \item If the observer measures again, then $S$ gets state $i$ with probability $1$ (nothing has changed).
\end{itemize}
\end{frame}

% --- Slide 37 ---
\begin{frame}{Classical physics}
\begin{itemize}
  \item If the system undergoes isolated evolution (i.e.\ “following the laws of physics”), then the state of the system gets updated via multiplication by a stochastic matrix $A$:
  \[
    \mathbf{p}\mapsto \mathbf{p}' = A\,\mathbf{p}.
  \]
  \item A matrix is \emph{stochastic} if entries are nonnegative, and each column sums to $1$.
  \item Stochastic matrices map probability vectors to probability vectors.
\end{itemize}
\end{frame}

% --- Slide 38 ---
\begin{frame}{Quantum physics}
\begin{itemize}
  \item Initially, the observer assigns a state to the system $S$.
  \item According to quantum physics, we can model the state of $S$ as a complex unit vector in $\mathbb{C}^d$, represented as a column vector:
  \[
    \lvert \psi\rangle =
    \begin{pmatrix}
      \alpha_0\\ \vdots\\ \alpha_{d-1}
    \end{pmatrix},\qquad
    \sum_{i=0}^{d-1}\lvert \alpha_i\rvert^2 = 1.
  \]
  \item The $d$ distinguishable states (also called “classical states”) are represented by the standard basis vectors
  \[
    \lvert 0\rangle =
    \begin{pmatrix}1\\0\\ \vdots\\0\end{pmatrix},\quad
    \lvert 1\rangle =
    \begin{pmatrix}0\\1\\ \vdots\\0\end{pmatrix},\quad \dots,\quad
    \lvert d-1\rangle =
    \begin{pmatrix}0\\0\\ \vdots\\1\end{pmatrix}.
  \]
  \item A general quantum state is a superposition of classical basis states:
  \[
    \lvert \psi\rangle = \alpha_0\lvert 0\rangle + \alpha_1\lvert 1\rangle + \cdots + \alpha_{d-1}\lvert d-1\rangle.
  \]
  \item The $\alpha_i$’s are called \emph{amplitudes}.
\end{itemize}
\end{frame}

% --- Slide 39 ---
\begin{frame}{Dirac notation}
\begin{itemize}
  \item The notation $\lvert \psi\rangle$ is called \emph{Dirac notation}, used to represent quantum states.
  \item Mathematically, a “ket vector” $\lvert \psi\rangle$ is a column vector.
  \item The dual/Hermitian conjugate of column vectors (i.e.\ row vectors) are called “bra vectors”:
  \[
    \langle \psi \rvert = (\lvert \psi\rangle)^\dagger.
  \]
  \item Example (qubit basis):
  \[
    \langle 0\rvert = (1,0),\qquad \langle 1\rvert = (0,1).
  \]
  \item Inner product notation:
  \[
    \langle \phi \mid \psi\rangle \in \mathbb{C}.
  \]
  \item Naming: “bra” + “ket” = “bracket”.
\end{itemize}
\end{frame}

% --- Slide 40 ---
\begin{frame}{Quantum physics}
\begin{itemize}
  \item If the observer measures the system $S$, then obtains a measurement outcome $i$ with probability
  \[
    \Pr[i] = \lvert \alpha_i\rvert^2,
  \]
  and then the state of the system gets updated (“collapsed”) from $\lvert \psi\rangle$ to the classical state $\lvert i\rangle$.
  \item If the observer measures again, then it gets state $i$ with probability $1$.
  \item This is called the \textbf{Born Rule}.
\end{itemize}
\end{frame}

% --- Slide 41 ---
\begin{frame}{Quantum physics}
\begin{itemize}
  \item If the system undergoes isolated evolution (i.e.\ “following the laws of physics”), then the state gets updated via multiplication by a unitary matrix $U$:
  \[
    \lvert \psi\rangle \mapsto \lvert \psi'\rangle = U\lvert \psi\rangle.
  \]
  \item A complex matrix $U$ is \emph{unitary} if
  \[
    U^\dagger U = I \qquad (\text{equivalently } U^{-1}=U^\dagger).
  \]
  \item $U^\dagger$ denotes the Hermitian conjugate: transpose then complex-conjugate every entry.
  \item Unitary matrices map unit vectors to unit vectors.
\end{itemize}
\end{frame}

% --- Slide 42 ---
\begin{frame}{Measuring a qubit}
\begin{itemize}
  \item Example state:
  \[
    \lvert \psi\rangle = \sqrt{\frac{2}{3}}\lvert 0\rangle + \sqrt{\frac{1}{3}}\lvert 1\rangle.
  \]
  \item After measurement:
  \begin{itemize}
    \item Outcome $0$ with probability $\frac{2}{3}$, post-measurement state $\lvert 0\rangle$.
    \item Outcome $1$ with probability $\frac{1}{3}$, post-measurement state $\lvert 1\rangle$.
  \end{itemize}
\end{itemize}
\end{frame}

% --- Slide 43 ---
\begin{frame}{Evolution of a qubit}
\begin{itemize}
  \item Example: start with $\lvert \psi\rangle = \lvert 1\rangle$.
  \item Apply a unitary $U$ (as shown on the slide) to get
  \[
    \lvert \psi'\rangle = -\frac{1}{\sqrt{2}}\lvert 0\rangle + \frac{1}{\sqrt{2}}\lvert 1\rangle.
  \]
\end{itemize}
\figplaceholder{Matrix $U$ shown on slide}
\end{frame}


\begin{frame}[t]{A demonstration of quantum weirdness}
Experiment B
\vspace{0.5em}
\begin{itemize}
  \item Transition
  \item amplitudes
\end{itemize}
\vspace{0.5em}
\end{frame}

\begin{frame}[t]{A demonstration of quantum weirdness}
\begin{itemize}
  \item Takeaway: Intermediate measurements
  \item destroy superpositions and prevent interference!
  \item Measurements make qubits become classical.
  \item Takeaway: to see quantum effects, delay
  \item measurements for as long as possible.
  \item This is why it is hard to see quantum effects
  \item in everyday life. Measurement is constantly
  \item happening.
\end{itemize}
\vspace{0.5em}
Experiment A
\vspace{0.5em}
\begin{itemize}
  \item Intermediate
  \item measurement
\end{itemize}
\vspace{0.5em}
Intermediate measurements prevent interference.
\vspace{0.5em}
\end{frame}

\begin{frame}[t]{What is a quantum state describing?}
\begin{itemize}
  \item A classical state is a probability distribution describing our uncertainty about the “true” underlying state of system S.
  \item When the observer measures, the “true” state is revealed. However, intermediate measurements affect outcomes.
  \item Quantum states are different: intermediate measurements do alter the system and cause different outcomes than if you didn’t measure.
  \item Thus quantum states are not describing uncertainty of a system. There is no “true” underlying classical state.
  \item But as external observers, we can only see ”classical shadows” of quantum states through measurement.
\end{itemize}
\vspace{0.5em}
\end{frame}




\end{document}
